\chapter{World Powers and Obligation}
\label{ch:world-powers-obligation}

\index{terrestrial patrons}\index{obligation}\index{factions}
Not every source of power wears a crown of stars. The guildmaster who controls the docks, the spymaster who knows your secret, the moneylender who holds your marker—these are patrons of a different kind. They grant no miracles, but they offer \textbf{leverage}: sanctuary, soldiers, silver, and silence.

This chapter covers \textbf{Terrestrial Patrons}—mortal factions and powerful individuals who can aid you, but at a price. The mechanics mirror those of supernatural Patrons (\S\ref{ch:magic}), but the consequences are social, legal, and economic rather than mystical.

\section{Terrestrial vs. Divine Patrons}

\begin{center}
\begin{tabular}{p{4cm}p{5cm}p{5cm}}
\toprule
& \textbf{Terrestrial Patrons} & \textbf{Divine Patrons} \\
\midrule
\textbf{Nature} & Mortal factions, guilds, nobles & Gods, numina, cosmic offices \\
\textbf{Benefits} & Resources, sanctuary, information, military aid & Miracles, rites, supernatural power \\
\textbf{Weight Track} & Same capacity (Spirit+Presence) & Same capacity \\
\textbf{Overage} & Fatigue (social/legal friction) & Fatigue (spiritual strain) \\
\textbf{Fallout} & Bounty, blacklist, legal trouble, rivals & Corruption, Patron intrusion, curses \\
\textbf{Clearing} & Service, bribes, political favors & Rituals, sacrifices, Patron tasks \\
\bottomrule
\end{tabular}
\end{center}

\emph{Both can lift you up. Both can weigh you down. One leaves a mark on your soul; the other leaves a mark on your reputation.}

\section{The Weight of Obligation (Refresher)}

Your \textbf{Obligation Capacity} is \textbf{Spirit + Presence}. Each time you call on a Patron’s aid (or accept their “gift”), you add one or more segments of \textit{claim} to that Patron.

\begin{itemize}
  \item \textbf{Exceeding Capacity:} Every segment over capacity inflicts 1 Fatigue (stress, social pressure, unwanted attention).
  \item \textbf{Clearing the Track:} Reduce your burden by performing service, paying bribes, completing tasks, or using downtime actions.
  \item \textbf{Full Clock:} At 6 segments (or when the GM decides), the Patron makes a \textbf{demand} you cannot easily refuse.
\end{itemize}

See \S\ref{ch:magic} for full Obligation mechanics. The same rules apply to terrestrial Patrons, but the \textbf{backlash} is always mundane: a rival faction is tipped off, a favor is called in early, or a former ally turns hostile.

\section{Major Terrestrial Patrons}

Below are examples of mortal powers that characters might serve (or drift into service with). Each entry includes:
\begin{itemize}
  \item \textbf{Typical Perks} – what you gain by spending a Boon or calling in a favor
  \item \textbf{Sample Demands} – what they may ask in return
  \item \textbf{Obligation Triggers} – when you mark +Obligation
  \item \textbf{Favor} – what a monumental service earns you
\end{itemize}

\subsection{The Thieves’ Guild (Syndicate)}

\textit{Every city has one. They own the shadows, the fences, and the watch’s blind spots.}

\textbf{Typical Perks:}
\begin{itemize}
  \item Sanctuary in a hidden safehouse for 1–3 nights.
  \item A forged document (identity, travel pass, customs waiver) usable once.
  \item Information about a target’s routines or vulnerabilities.
  \item A “clean” fence for stolen goods (no questions, lower return).
\end{itemize}

\textbf{Sample Demands:}
\begin{itemize}
  \item “Look the other way” during a heist you witness.
  \item Smuggle a small package across a border.
  \item Plant evidence on a rival.
  \item Pay a “protection fee” from your own pockets.
  \item Recruit a new member or identify a promising mark.
\end{itemize}

\textbf{Obligation Triggers:}
\begin{itemize}
  \item Using any perk: +1 weight.
  \item Requesting a major heist or sanctuary during a manhunt: +2.
  \item Refusing a demand outright: +1 (and possible retaliation).
\end{itemize}

\textbf{Favor (reduce burden by 2, gain permanent asset):}
\begin{itemize}
  \item \textbf{Guild Status:} You are marked as a trusted associate; fences offer better rates.
  \item \textbf{One “Get Out of Jail” card:} The Guild will break you out of a non‑magical prison once.
\end{itemize}

\subsection{Mercenary Company (Condotta)}

\textit{Banners for hire. They fight for coin, but loyalty can be bought with more than silver.}

\textbf{Typical Perks:}
\begin{itemize}
  \item A squad of 4–6 soldiers (Cap 2 each) for one battle or scene.
  \item Use of a fortified camp or watchtower as a base.
  \item Military intelligence about a region (troop movements, supply lines).
  \item A captain’s letter of marque (legal cover for violence against enemies of the company).
\end{itemize}

\textbf{Sample Demands:}
\begin{itemize}
  \item Fight alongside them in a skirmish (no extra pay).
  \item Forge a payroll document or help loot a supply wagon.
  \item Sponsor a new recruit or vouch for a deserter.
  \item Carry a message to a rival company (may be a trap).
\end{itemize}

\textbf{Obligation Triggers:}
\begin{itemize}
  \item Using any perk: +1.
  \item Requesting a major deployment (20+ soldiers): +2.
  \item Deserting in battle or refusing a direct order: +2 and possible bounty.
\end{itemize}

\textbf{Favor:}
\begin{itemize}
  \item \textbf{Honorary Rank:} You command a small detachment (Cap 3 followers) for the arc.
  \item \textbf{Share of Plunder:} A permanent +1 to resource rolls when working with the company.
\end{itemize}

\subsection{Bank or Charter House (Coin-Prince)}

\textit{They hold the realm’s debts. Their ledgers are law, and their vaults are legend.}

\textbf{Typical Perks:}
\begin{itemize}
  \item A letter of credit for 100 gold (or equivalent liquid assets).
  \item A day’s grace on a due payment (postpone a financial complication).
  \item Access to a private vault or secure courier.
  \item An introduction to a noble factor or merchant consortium.
\end{itemize}

\textbf{Sample Demands:}
\begin{itemize}
  \item Repay a loan with interest (narrative timeframe).
  \item Act as a debt collector (intimidation or persuasion vs. a debtor).
  \item Transport a sealed chest (no questions) across a border.
  \item Invest in a failing enterprise (lose XP or suffer a financial setback).
\end{itemize}

\textbf{Obligation Triggers:}
\begin{itemize}
  \item Using any perk: +1.
  \item Large loan (200+ gold) or co‑signing a risky venture: +2.
  \item Defaulting on repayment: track fills, and the Bank seizes assets (lose a Minor Asset).
\end{itemize}

\textbf{Favor:}
\begin{itemize}
  \item \textbf{Preferred Status:} No interest on future loans; +1 die on commerce rolls in that city.
  \item \textbf{Minor Share:} A permanent small income (enough to cover one follower’s upkeep).
\end{itemize}

\subsection{Royalist Clandestine Network (The Sun’s Shadow)}

\textit{Agents of a deposed or contested crown. They work in whispers, seeking to restore their monarch.}

\textbf{Typical Perks:}
\begin{itemize}
  \item A safe house in a loyalist neighborhood.
  \item A coded message that opens doors among certain nobles.
  \item A forged pedigree or passport (usable once).
  \item Warning of an impending raid or arrest warrant.
\end{itemize}

\textbf{Sample Demands:}
\begin{itemize}
  \item Deliver a sealed letter to a “sleeping” agent.
  \item Create a diversion during a loyalist meeting.
  \item Smuggle a fugitive noble across the border.
  \item Plant a rumor that discredits a rival claimant.
\end{itemize}

\textbf{Obligation Triggers:}
\begin{itemize}
  \item Using any perk: +1.
  \item Requesting a major operation (rescuing a captured agent): +2.
  \item Betraying the network: Permanent enemy faction; track becomes a Bounty clock.
\end{itemize}

\textbf{Favor:}
\begin{itemize}
  \item \textbf{Title Restored:} If the claimant regains power, you receive a minor title or land grant.
  \item \textbf{Secret Keeper:} You learn a dangerous secret about a rival house—blackmail material.
\end{itemize}

\subsection{Ecktorian Mystery Cult (The Unseen Flame)}

\textit{These cults hide within the Everflame’s temples. They seek hidden knowledge, forbidden rites, and influence over the living flame.}

\textbf{Typical Perks:}
\begin{itemize}
  \item One free use of a temple’s sacred flame (purification, healing, or divination effect).
  \item A cipher key to read encoded Ecktorian legal documents.
  \item A minor relic that grants +1 die to resist fear or fire.
  \item A secret passage within the city’s catacombs.
\end{itemize}

\textbf{Sample Demands:}
\begin{itemize}
  \item Retrieve a heretical text from a rival sect.
  \item Silence a whistleblower who threatens the cult’s secrecy.
  \item Perform a ritual at a specific hour (may attract attention).
  \item Provide a new member (a vulnerable person to indoctrinate).
\end{itemize}

\textbf{Obligation Triggers:}
\begin{itemize}
  \item Using any perk: +1.
  \item Learning a major secret (e.g., a high priest’s corruption): +2.
  \item Refusing a direct order from a cult leader: +2 and possible blackmail.
\end{itemize}

\textbf{Favor:}
\begin{itemize}
  \item \textbf{Initiate’s Mark:} You can invoke the cult’s sigil to lower DV of one social roll with Everflame clergy.
  \item \textbf{Blessed Flame:} You may light a sacred candle once per session to gain a free Boon.
\end{itemize}

\subsection{Imperial Remnant Organization (The Seventh Legion)}

\textit{Remnants of the old Utaran Empire—veterans, archivists, and dreamers who believe the empire can be reborn.}

\textbf{Typical Perks:}
\begin{itemize}
  \item Access to an imperial archive (useful for Lore or Investigation rolls).
  \item A map of old military roads (faster travel in certain regions).
  \item A “coin of passage” recognized by old legion outposts.
  \item A retired centurion (Cap 2 follower) for one mission.
\end{itemize}

\textbf{Sample Demands:}
\begin{itemize}
  \item Recover a lost legion banner from a haunted battlefield.
  \item Sabotage a rival’s attempt to claim imperial salvage.
  \item Recruit a new member with military or political skills.
  \item Publicly renounce a rival faction (e.g., the Black Banners).
\end{itemize}

\textbf{Obligation Triggers:}
\begin{itemize}
  \item Using any perk: +1.
  \item Requesting a small cohort (10–20 veterans): +2.
  \item Betraying the remnant: Permanent enemy faction; track becomes a Vendetta clock.
\end{itemize}

\textbf{Favor:}
\begin{itemize}
  \item \textbf{Legionary Title:} You are recognized as a \emph{Tribune} – legal authority in old imperial territories.
  \item \textbf{Stored Cache:} Access to a hidden armory (Minor Asset: weapons and armor).
\end{itemize}

\section{Creating Your Own Terrestrial Patrons}

Use these steps to build a custom faction:

\begin{enumerate}
  \item \textbf{Nature:} What kind of organization? (Thieves, merchants, clergy, military, noble house, etc.)
  \item \textbf{Resources:} What can they realistically provide? (Perks – 2 to 4 concrete benefits)
  \item \textbf{Agenda:} What do they want? (Demands – 3 to 5 typical tasks)
  \item \textbf{Price:} What does taking their help cost? (Obligation triggers)
  \item \textbf{Reward:} What does truly exceptional service earn? (Favor – 1 major permanent boon)
\end{enumerate}

\paragraph{Example: The Dockmen’s Union (Custom)}  
\textit{Nature:} Labor guild controlling river freight.  
\textit{Perks:} Priority berthing, free loading crew, rumor mill about customs schedules.  
\textit{Demands:} Break a strike, intimidate a rival captain, look the other way during a smuggling run.  
\textit{Obligation Triggers:} Using a perk = +1; requesting armed protection = +2.  
\textit{Favor:} A permanent shipping contract (income + trade route intel).

\section{Managing Multiple Terrestrial Patrons}

Characters may serve (or owe) more than one terrestrial patron. Keep separate Obligation tracks for each. When two patrons make conflicting demands, the GM may force a choice – refusing one adds +2 to that track and may trigger immediate consequences (e.g., the other patron gains a Boon to use against you).

\emph{Example:} You carry a marker for the Thieves’ Guild (3) and a Royalist network (2). The Guild demands you rob a royal courier; the Royalists ask you to protect that same courier. You cannot satisfy both. Choose one: the other track fills by 2, and the spurned patron becomes a Rival.

\section{Terrestrial Obligation and the Game World}

Unlike divine Obligation, which often manifests as personal corruption or Patron intrusions, terrestrial Obligation creates \textbf{worldly pressure}:
\begin{itemize}
  \item \textbf{Gossip and reputation:} Your name appears on blacklists.
  \item \textbf{Legal trouble:} Fines, arrests, or revoked licenses.
  \item \textbf{Economic pressure:} Prices rise, trade opportunities vanish.
  \item \textbf{Violence:} Rivals send thugs, assassins, or challenge you to duels.
  \item \textbf{Betrayal:} A former ally sells you out to another faction.
\end{itemize}

The GM should make these consequences visible and actionable – a chance to pay off outstanding favors, shift allegiances, or break free entirely.

\section{Optional Rule: Faction Clocks}

For a longer campaign, track each terrestrial patron’s \textbf{Faction Clock} (size 6–8). The clock advances when:
\begin{itemize}
  \item You call in a major favor (+1)
  \item You refuse a demand (+1)
  \item You fail a mission (+2)
  \item You work against their interests (+2)
\end{itemize}

When the clock fills, the patron makes a \textbf{power play}: they might seize an asset, harm an ally, reveal a secret, or demand an impossible task. The clock resets, but the relationship is permanently strained.

\begin{tcolorbox}[colback=gray!5!white,colframe=gray!75!black,title=World Powers Quick Reference,fonttitle=\bfseries]
\textbf{Obligation Capacity:} Spirit + Presence \\
\textbf{Over Capacity:} Fatigue per segment; at 6 segments = demand \\
\textbf{Typical Terrestrial Patrons:} Thieves’ Guild, Mercenary Company, Bank, Royalist Network, Mystery Cult, Imperial Remnant \\
\textbf{Common Perks:} Sanctuary, forged documents, information, soldiers, loans \\
\textbf{Common Demands:} Service, silence, sabotage, recruitment \\
\textbf{Favor Reward:} Permanent title, asset, or reputation bonus
\end{tcolorbox}

\begin{tcolorbox}[colback=black!3!white,colframe=black!60!white,title={From Faction to Asset: Taking Over a Terrestrial Patron}]
  \index{terrestrial patrons!takeover}\index{assets!from factions}
  Sometimes you stop owing the guild – and start owning it. Rising to lead a faction is a major story arc, not a simple XP purchase. This sidebar outlines how to turn a terrestrial patron into a personal Asset.
  
  \paragraph{Narrative Prerequisites}
  Before you can attempt a takeover, you must:
  \begin{itemize}
    \item Have at least \textbf{6 Obligation} with the faction (you are deeply entangled).
    \item Complete a \textbf{major service} for the faction (GM’s discretion – e.g., saving it from destruction, winning a war, exposing a traitor at great cost).
    \item Secure at least one \textbf{internal ally} with significant influence (another captain, a guildmaster, a wealthy backer).
  \end{itemize}
  
  \paragraph{The Takeover Arc}
  The GM and player collaborate on a 2–3 session arc that typically involves:
  \begin{enumerate}
    \item \textbf{Crisis or Power Vacuum} – The current leader dies, disappears, or is disgraced.
    \item \textbf{Claim of Right} – You present a valid claim (charter, election, duel victory, popular support).
    \item \textbf{Consolidation} – Defeat or co‑opt rivals, bribe or intimidate dissenters, and secure the faction’s resources.
  \end{enumerate}
  
  \paragraph{Mechanical Cost}
  Once the arc is complete, you must spend \textbf{12 XP} (the cost of a Major Asset) to convert the faction into a personal Asset. This represents the time, bribes, and institutional restructuring required to make it \emph{yours} rather than a temporary alliance.
  
  \paragraph{Converting the Asset}
  The new Asset uses the statistics of the most appropriate existing Asset from this compendium, modified by the faction’s nature:
  \begin{itemize}
    \item \textbf{Thieves’ Guild} → \textbf{Informant Ring} (Standard) or \textbf{Smuggler’s Network} (Standard)
    \item \textbf{Mercenary Company} → \textbf{Mercenary Contract} (Minor/Standard) or \textbf{Military Outpost} (Major)
    \item \textbf{Banking Charter} → \textbf{Banking Charter} (Major) or \textbf{Trading Consortium} (Major)
    \item \textbf{Royalist Network} → \textbf{Intelligence Bureau} (Specialized) or \textbf{Diplomatic Enclave} (Major)
    \item \textbf{Ecktorian Mystery Cult} → \textbf{Temple Complex} (Standard) or \textbf{Magical Laboratory} (Specialized)
  \end{itemize}
  The GM may also create a \textbf{custom Asset} using the guidelines in \S\ref{sec:creating-assets}.
  
  \paragraph{What You Gain}
  \begin{itemize}
    \item The Asset’s normal \textbf{Free Use} and \textbf{Scene Surge} benefits.
    \item The \textbf{Access Tag} associated with the faction (e.g., \textit{Guild-Ledgered}, \textit{Harbor Handsigns}).
    \item A permanent \textbf{obligation track} toward the faction’s members (now your underlings). This track is smaller (4 segments) and fills when you neglect their welfare or make reckless decisions that endanger them.
  \end{itemize}
  
  \paragraph{What You Lose}
  \begin{itemize}
    \item You can no longer use the faction as a \emph{terrestrial patron} to gain Obligation‑free perks – you are now responsible for it.
    \item The faction becomes a target for your enemies. The GM may advance a \textbf{Rival Faction clock} whenever you use the Asset aggressively.
    \item You must pay normal \textbf{upkeep} (XP or downtime scenes) to keep the Asset \emph{Maintained}. Neglect causes internal dissent, corruption, or defections.
  \end{itemize}
  
  \paragraph{Optional: Gradual Takeover (No XP Cost)}
  For a slower, more story‑driven approach, the player can forgo the 12 XP cost but must instead \textbf{spend a downtime scene per session} managing the faction (equivalent to Intensive upkeep) and accept that the Asset’s benefits are unreliable until they have successfully completed a full arc of consolidation. The GM may allow the XP cost to be paid in installments (e.g., 4 XP per downtime period) to represent gradual investment.
  
  \paragraph{Example: Rellan Takes the Black Ledger}
  Rellan has served the Black Ledger smuggling syndicate for two arcs, accruing 6 Obligation. After exposing a rival who betrayed the syndicate, the aging leader steps down. Rellan spends 12 XP, completes a short arc consolidating loyal captains, and converts the syndicate into a \textbf{Smuggler’s Network (Standard Asset)}. He now commands hidden routes and safehouses – but must keep the crew paid and protected, or an internal revolt clock will start ticking.
  \end{tcolorbox}
  
\section{Narrative Hooks by Region}

\subsection*{Silkstrand}
\begin{itemize}
  \item \textbf{The Dyers’ Cartel:} They control the dye trade – and the curses that stain traitors’ cloth.
  \item \textbf{The Condotta Registry:} A broker for mercenary contracts. They know where every sword is pointed.
\end{itemize}

\subsection*{Ecktoria}
\begin{itemize}
  \item \textbf{The Everflame Audit:} A secret committee of Lampers who keep a black book of corrupt nobles.
  \item \textbf{The Marble Forum:} A law society that can stall or accelerate any legal proceeding – for a price.
\end{itemize}

\subsection*{Kahfagia}
\begin{itemize}
  \item \textbf{The Lantern-Law Convoys:} Merchant captains who share beacon codes. Join them and you sail faster – but you must leave a family member as hostage.
\end{itemize}

\subsection*{Mistlands}
\begin{itemize}
  \item \textbf{The Bell-Wardens:} They can silence your name in the bells – or toll your death.
  \item \textbf{The Salt-Circle:} An apothecary guild that also trades in memories harvested from the fog.
\end{itemize}

\subsection*{Vhasia / Viterra}
\begin{itemize}
  \item \textbf{The Sun’s Shadow (royalist network)} – see above.
  \item \textbf{The Hedge-Law Clerks:} They specialize in “adjusting” property boundaries and inheritance documents.
\end{itemize}

\section{Closing Advice}

Terrestrial patrons make the political world feel alive. They offer immediate, practical aid – but every coin comes with a string. Use them to create moral dilemmas, escalate rivalries, and reward long-term investment in the setting.

\emph{Remember: The gods demand your soul. The guildmaster only wants your silence – or a small favor, delivered by midnight.}